\section{Ejercicio 1: entrenamiento comparativo (CBOW frente a Skip-gram)}

\pregunta{1.1}{(0,5 punto) ¿Qué diferencias observan en las palabras más cercanas retornadas por cada arquitectura? Explica en forma concisa y precisa los aspectos destacables de la diferencia.}
\textbf{Respuesta 1.1.} En la práctica las dos arquitecturas devuelven casi los mismos vecinos para \code{cobre} (\code{solar} y \code{local} salen en ambas listas), así que la diferencia de fondo no está en las palabras sino en los puntajes: CBOW da similitudes entre 0,62 y 0,64, mientras que en Skip-gram los tres primeros superan 0,91 y quedan casi empatados. Esto se explica por la forma de entrenar de cada uno. Skip-gram genera un ejemplo por cada par (palabra, contexto), y con 200 épocas sobre un corpus de apenas 12 oraciones el espacio termina muy comprimido: casi todos los vectores apuntan en direcciones parecidas, las similitudes se inflan y el modelo pierde capacidad de discriminar. CBOW en cambio promedia el contexto antes de predecir, lo que suaviza las actualizaciones y deja un espacio menos saturado, donde los puntajes todavía sirven para distinguir un vecino de otro. Nos llamó la atención además que ninguno de los dos entregue palabras ``mineras'' para \code{cobre}, pero revisando el corpus tiene sentido: \code{cobre} aparece solo en dos oraciones (una sobre eficiencia energética y otra sobre ingresos fiscales), y con tan poca evidencia el modelo no alcanza a anclarlo al dominio minero; sus vecinos terminan siendo palabras con las que simplemente coocurre seguido.

\pregunta{1.2}{(1 punto) Modifica el parámetro window entre 3 y 10. ¿Qué sucede con las palabras más cercanas a 'cobre'? ¿Se vuelven más especializadas (relacionadas con minería) o más genéricas (relacionadas con la estructura de la oración)? Justifica y detalla de forma concisa y precisa.}
\textbf{Respuesta 1.2.} Al probar con ventanas de 3, 5, 7 y 10 (celda anterior), los vecinos de \code{cobre} no se especializan en minería; pasa lo contrario: se vuelven más genéricos y cada vez más difíciles de separar entre sí. Con \code{window=10} en Skip-gram los tres primeros quedan entre 0,980 y 0,982, es decir, la geometría se aplana tanto que casi cualquier palabra del corpus queda ``cerca'' de \code{cobre}. En teoría una ventana chica captura relaciones más locales y sintácticas (ligadas a la posición de la palabra en la oración) y una grande captura relaciones temáticas, pero acá pesa más un tema práctico: después de sacar las stopwords cada oración queda con 7 a 9 tokens, así que una ventana de 10 cubre la oración completa y el contexto de cada palabra pasa a ser todo su documento. En esas condiciones dos palabras terminan siendo similares solo por compartir oración, y eso homogeneiza el espacio. De hecho los vecinos que se repiten en todas las configuraciones (\code{requiere}, \code{local}, \code{solar}) son palabras que coocurren con \code{cobre} o que aparecen en varias oraciones, no términos propios de la minería. En un corpus de este tamaño agrandar la ventana no especializa nada, solo diluye la señal.

\section{Ejercicio 2: álgebra vectorial y analogías semánticas}

\pregunta{2.1}{(0,5 punto) Interpretación del Espacio Vectorial: Dado que el resultado de la operación v\_mineria - v\_cobre + v\_inflacion arrojó como palabras más cercanas 'tasa' y 'banco', ¿qué relación semántica está capturando el modelo sobre la 'Inflación' en comparación con la relación entre 'Minería' y 'Cobre'?}
\textbf{Respuesta 2.1.} Un detalle antes de interpretar: como a \code{most\_similar} le pasamos el vector crudo y no los nombres de las palabras, la función no excluye los términos de entrada, y por eso \code{minería} (0,8209) e \code{inflación} (0,7978) encabezan la lista (el vector resultante conserva buena parte de ambos, así que era de esperar). Lo que vale la pena mirar son las primeras palabras nuevas: \code{impacto} (0,7432), \code{tasa} (0,7218) y \code{crediticio} (0,7094), todas del ámbito financiero-monetario. La relación que está capturando el modelo es la de un concepto con el dominio que lo gestiona: igual que \code{cobre} es el objeto material de la \code{minería}, \code{inflación} es la variable central del aparato monetario, donde la \code{tasa} es el instrumento para controlarla y el \code{banco} (central) la institución que la administra. La resta \code{v\_minería - v\_cobre} aísla, aunque sea de forma gruesa, la dirección de ``actividad que engloba al objeto'', y al sumarla a \code{inflación} el vector cae en la zona del espacio donde están los términos de política monetaria. Que esto funcione con 12 oraciones se debe a que el corpus es muy temático: \code{inflación} coocurre únicamente con \code{banco}, \code{central} y \code{tasa}.

\pregunta{2.2}{(0,5 punto) Si el modelo no cuenta con suficientes datos para aprender la relación 'Cobre-Minería', la analogía falla. Menciona una estrategia (ej. aumentar corpus, usar embeddings pre-entrenados como word2vec-google-news-300) para mejorar la precisión de estas analogías.}
\textbf{Respuesta 2.2.} Lo que más rinde es partir de embeddings pre-entrenados sobre corpus masivos: \code{word2vec-google-news-300} si se trabaja en inglés o, para nuestro caso, alternativas en español como los vectores del Spanish Billion Word Corpus o fastText en español. Esos modelos ya vieron millones de contextos donde \code{cobre} aparece junto a \code{minería}, \code{yacimiento} o \code{Codelco}, así que la dirección de la analogía viene aprendida de manera robusta, mientras que en nuestro corpus esa relación cuelga de apenas dos coocurrencias. La otra vía es ampliar el corpus, por ejemplo con miles de noticias económicas chilenas; también funciona, pero llegar al mismo nivel cuesta mucho más. Entre las opciones pre-entrenadas, fastText tiene la gracia de que trabaja con subpalabras y por eso generaliza incluso a términos que nunca vio en entrenamiento.

\pregunta{2.3}{(1 punto) Sensibilidad de la Analogía: Si cambiáramos el modelo de Skip-gram a CBOW para realizar esta misma operación aritmética, ¿esperarías que el vector resultante (v\_target) se acerque a los mismos conceptos? Justifica tu respuesta considerando cómo cada arquitectura distribuye los pesos semánticos en el vocabulario.}
\textbf{Respuesta 2.3.} Repetimos la operación con CBOW (celda anterior) y la dirección semántica se mantiene solo en parte: \code{impacto} y \code{crediticio} reaparecen, señal de que el vector sigue apuntando a la zona financiera, pero la analogía se degrada bastante. Las similitudes caen a menos de la mitad (0,4038 contra 0,7432 para la primera palabra nueva) y \code{tasa}, que era el resultado más limpio en Skip-gram, desaparece del top-5; en su lugar entra \code{influyen}, que no dice mucho. El porqué está en cómo reparte los pesos cada arquitectura: CBOW promedia los vectores del contexto antes de predecir, lo que suaviza las representaciones y hace que las palabras frecuentes dominen el aprendizaje, mientras que las poco frecuentes (y acá casi todo el vocabulario aparece una o dos veces) reciben actualizaciones diluidas. Skip-gram entrena cada par palabra-contexto por separado y deja a los términos raros con vectores mejor definidos. Como la aritmética de analogías necesita que \code{cobre}, \code{minería} e \code{inflación} tengan direcciones nítidas, la arquitectura que mejor representa palabras infrecuentes, o sea Skip-gram, entrega el resultado más interpretable: con CBOW el concepto general se conserva, pero la precisión se pierde.

\section{Ejercicio 3: reducción de dimensionalidad y visualización}

\pregunta{3.1}{Análisis de Proyección PCA Al observar el gráfico de 2D, ¿identificas alguna agrupación lógica de palabras (clusters) que refleje las temáticas del corpus (Minería, Energía, Finanzas)? Explica si la cercanía espacial entre términos como 'cobre' y 'minería' o 'inflación' y 'tasa' se mantiene tras reducir la dimensionalidad de 60 a 2 dimensiones.}
\textbf{Respuesta 3.1.} Antes de buscar clusters conviene tener presente cuánta información sobrevive a la proyección: entre PC1 y PC2 se explica apenas un 24,4 \% de la varianza (19,8 \% y 4,6 \%), así que el gráfico es un resumen bastante parcial del espacio de 60 dimensiones. Aun así, sí se distinguen agrupaciones lógicas, aunque más como vecindades difusas que como clusters nítidos. La más clara es la financiera: \code{inflación}, \code{banco}, \code{central}, \code{interés}, \code{economía}, \code{evalúan} y \code{controlar} quedan en una franja contigua en la zona centro-inferior del plano. Arriba, hacia el centro-derecha, se arma una vecindad minero-energética con \code{minería}, \code{verde}, \code{agua}, \code{litio}, \code{energía}, \code{eólica} y \code{evaluaciones}. Sobre las cercanías puntuales, el resultado es desigual (los números están en la celda anterior). \code{inflación}-\code{tasa}, con coseno 0,872 en 60D, queda a distancia 0,032 en 2D, y \code{banco}-\code{central} (0,825) a 0,017; las dos por debajo de la distancia media del gráfico (0,054), o sea que esas relaciones sobreviven a la proyección. \code{cobre}-\code{minería}, en cambio, tiene una similitud altísima en 60D (0,876) pero queda a 0,050 en 2D, prácticamente la distancia promedio, y de hecho en el gráfico \code{cobre} aparece en el extremo derecho, lejos de \code{minería}. Ahí se ve la limitación de PCA: al ser una proyección lineal que maximiza la varianza global, no garantiza conservar las vecindades locales, y la dirección que unía \code{cobre} con \code{minería} simplemente no quedó alineada con los dos primeros componentes. Para estudiar vecindades locales convendría t-SNE, que optimiza justamente eso en vez de la varianza global.

\section{Ejercicio 4: modelado de tópicos con document embeddings}

\pregunta{4.1}{(0,8 punto) \textbf{Calidad del Agrupamiento (Coherencia)}: Al observar los documentos asignados a cada cluster, ¿consideras que el modelo logró separar adecuadamente las temáticas de 'Minería/Energía' de las de 'Finanzas/Economía'? Identifica si hay algún documento que parezca 'fuera de lugar' y explica por qué el promedio de vectores (Mean Pooling) podría haber causado esa asignación.}
\textbf{Respuesta 4.1.} La separación funcionó a medias. El cluster 3 quedó impecable: aisló justo los dos documentos de política monetaria, que comparten casi el mismo vocabulario (\code{banco}, \code{central}, \code{tasa}, \code{inflación}). El cluster 2 también se puede interpretar, con un matiz: junta los dos documentos de litio con el del precio del cobre, así que más que minería pura arma un tópico de ``commodities y su valor económico''. El problema es el cluster 1, que quedó como cajón de sastre: mezcla documentos minero-energéticos con tres claramente financieros (``instituciones financieras evalúan el riesgo crediticio'', ``sistema financiero mantiene estándares de liquidez'', ``analistas proyectan crecimiento del PIB'') que por tema deberían haber caído junto al cluster 3. Esos tres son los que están fuera de lugar, y el mean pooling tiene harto que ver: al promediar todos los vectores del documento, los términos discriminativos (\code{crediticio}, \code{liquidez}) se diluyen entre palabras genéricas que cruzan todo el corpus (\code{local}, \code{requiere}, \code{proyectan}, \code{evalúan}), que como vimos en el Ejercicio 1 son justamente las que el modelo dejó cerca de todo. El promedio arrastra los document embeddings hacia el centro común del espacio, y esos documentos financieros terminan más cerca del centroide ``genérico'' del cluster 1 que del núcleo monetario, que es chico y compacto.

\pregunta{4.2}{(0,7 punto) \textbf{Representación Global vs. Local}: ¿Qué sucede con palabras muy frecuentes pero poco informativas (como 'la', 'el', 'en') al calcular el promedio del documento? Propón una alternativa o mejora técnica al simple promedio para que los términos clave de la industria tengan mayor peso en la representación final del documento.}
\textbf{Respuesta 4.2.} En este taller los artículos y preposiciones se filtraron con las stopwords de NLTK antes de entrenar, así que no contaminan directamente el promedio. El fenómeno de fondo igual sigue presente, ahora con lo que podríamos llamar stopwords de dominio: palabras como \code{requiere}, \code{local} o \code{proyectos}, frecuentes en el corpus pero poco útiles para distinguir tópicos, pesan en el promedio exactamente igual que \code{litio} o \code{crediticio}. Si no se hubieran quitado las stopwords sería peor todavía: al ser las más frecuentes, dominarían el promedio y empujarían todos los documentos hacia un mismo punto del espacio, borrando las diferencias temáticas. Una mejora directa es ponderar el promedio con TF-IDF, multiplicando cada vector por la relevancia del término en el documento respecto del corpus: así \code{crediticio} (raro y discriminativo) pesaría mucho más que \code{requiere} (transversal). Una variante más fina es SIF (\textit{Smooth Inverse Frequency}, Arora et al., 2017), que pondera cada palabra por \code{a/(a + p(w))} y además remueve el primer componente principal común a todos los documentos. Y si se quiere dejar de lado el promedio, Doc2Vec o encoders de oraciones tipo SBERT aprenden la representación del documento directamente, sin armarla por agregación.
